\section{The cone tube and integral gluing}\label{sec:tube}

We now compute the integral middle homology and ordered Seifert
pairing for the functions constructed in the preceding section.
The calculation keeps the projection and transfer of each graph
map visible: these maps determine how cycles close across the
boundary of the smoothing region. Throughout this calculation
\(k=4\); the letter \(k\) records the covering degree in the formulas.

Let \(p=(X,T)\), \(h=T(X-T^2)(X-2T^2)\), and let
\(\pi:(\C^2,0)\to(\C^2,0)\) be one of the two degree-four finite
graph maps \(\pi_i\), with \(i\in\{0,1\}\) in the construction. Write \(d=h\circ\pi\), and let
\[
 R(p,w)=w^k+\sum_{j<k}A_j(p)w^j,\qquad k=4,
\]
be its prepared equation, so \(A_j(0)=0\) and \(R(0,w)=w^k\).
The root surface \(R=0\) is smooth near the origin; one component of \(R_p\)
is a unit. For \(s>0\), put \(R_s(p,w)=s^kR(p,w/s)\) and
\[
C_s=\{(p,y,z,w):yz=R_s(p,w)\},\qquad
f_s=h(p)+P(y,z,w),\quad P=w(y^a+z^b).
\]
The germ \(C_s\) is smooth of complex dimension four. The tube
statement below is formulated for any sufficiently large positive
integers \(a,b\) satisfying its two displayed inequalities. For the
two germs of this paper we use the single tuple already fixed in
Lemma~\ref{lem:simultaneous-parameters}. For such a qualifying pair, set
\[
\alpha=\frac{kb}{a+b},\qquad \beta=\frac{ka}{a+b},\qquad
\alpha+\beta=k,\qquad a\alpha=b\beta.
\]
Let \(\Hh,\D\) be the pages of \(h,d\). Let \(Q\) be a small
fiber \(P^{-1}(c)\) in the cone \(\{yz=w^k\}\), with \(c\ne0\);
it lies entirely in the smooth nonvertex locus.
Their middle homology groups are
\(H=H_1(\Hh;\Z)\), \(D=H_1(\D;\Z)\), and \(E=H_1(Q;\Z)\).
The finite map has projection and transfer \(\pi_*:D\to H\),
\(\pi^!:H\to D\), with \(\pi_*\pi^!=k\). Put
\(K=\ker\pi_*\) and
\(\pr l=l-k^{-1}\pi^!\pi_*l\in K_\Q\).
The cone link \(L\) is a lens space, \(H_1(L)=\Z/k\), and
\(\tau:E\to\Z/k\) is induced by weighted radial projection
\(Q\to L\), equivalently by page inclusion in the link open-book model.
Write \(s_h,s_d\) for the tube filling forms in the ordered
neighboring-page convention fixed in Section~\ref{sec:plane-cover}.
The cyclic map \((U,V)\mapsto(U^k,V^k,UV)\) to the cone pulls
\(P\) back to \(\widetilde P(U,V)=UV(U^{ka}+V^{kb})\).
Let \(\widetilde Q\) be its plane Milnor page, and let
\(\widetilde s\) be its plane-page form. The cyclic cover of the
nonvertex cone gives a full-lift map
\(\tr:E\to H_1(\widetilde Q;\Z)\); define
\[
 e(v,v')=\frac1k\widetilde s(\tr v,\tr v').
\]
This form is a priori rational because of the factor \(1/k\).
The theorem below specifies the integral page lattice on which
the rational form is evaluated.

\begin{theorem}[Tube, lattice, and pairing]\label{thm:tube}
For sufficiently large \(a,b\) satisfying
\[
 \frac{(a-2)\alpha}{k}>L_0,\qquad
 \frac{(b-2)\beta}{k}>L_0,
\]
where \(L_0=4\) is the gradient exponent in
\eqref{eq:plane-gradient}, and for sufficiently small \(s>0\), the germ
\(f_s:C_s\to\C\) has its sole critical point at the origin in a fixed
small tube. Its integral Milnor lattice, in rational coordinates, is
\begin{equation}
\Lambda_\pi=
\left\{(\pr l,z): l\in D,\ z\in H\otimes E,\quad
\pi_*l\bmod kH=(1\otimes\tau)z\right\}
\subset K_\Q\oplus(H\otimes E)_\Q.
\label{eq:tube-1}
\end{equation}
The rational orthogonal filling form on these coordinates is
\begin{equation}
  -s_d|_K\ \perp\ (s_h\otimes e).
\label{eq:tube-2}
\end{equation}
The lattice and form are those of a genuine local Milnor page of
\(f_s\) in the smooth ambient germ \(C_s\).
\end{theorem}

In \eqref{eq:tube-1}, \(z\) records the exterior product class,
while \(l\) indexes an inner closing chain. The equality of residues
is the condition for these pieces to close, and \(\pr l\) is the
resulting rational inner coordinate.

Section~\ref{sec:boundary-geometry} supplies boundary regularity,
compatible continuations, and isolation.
Section~\ref{sec:circle-collapse} computes the inner homology and
proves acyclicity of the large total. It then compares the large and
local total/page pairs, identifying the page homology and ordered
form with the ordinary local Milnor data. Its signed exterior cut
determines which exterior cycles extend.
Section~\ref{sec:integral-lattice}
then determines the integral lattice in these rational coordinates through the
Bockstein and transfer, and Section~\ref{sec:filling-summands}
computes the two ordered filling forms.
